\chapter{Reconstructing arguments with schemes and critical questions}
\label{sec:p18-argument-schemes-critical-questions}

An argument scheme identifies a kind of reason: an expert's judgment,
a witness's report, a comparison with another case, a proposed causal
connection, or a means to a goal. Identifying that reason helps an analyst
ask what the argument actually establishes and what could defeat its
support. A scheme is not a certificate that its conclusion is true.
Nor does an unanswered question by itself establish the opposite conclusion.

This chapter uses identified presentations by Godden and Walton
\cite{GoddenWalton2007SchemesQuestions}, Walton
\cite{Walton2014AnalogySchemes,Walton2016ArgumentEvaluation}, and
Macagno, Walton, and Reed \cite{MacagnoWaltonReed2017Schemes}.
The numbered question lists below preserve each selected source's
questions and distinctions in independently worded form.
Different lists are not combined into a supposedly universal checklist.
The worked cases are constructed examples, not reports of empirical
studies or accounts of actual people.

\section{A reconstruction and assessment procedure}

\begin{hbdefinition}{def:p18-scheme-assessment-terms}{Claim, warrant, qualifier, and challenge}
The \emph{conclusion} is the claim for which support is offered.
The \emph{premises} are the offered reasons.
A \emph{warrant} is the connection claimed between those reasons and
that conclusion; it may be implicit.
A \emph{qualifier} limits the force or scope of the conclusion:
for example, probably, normally, on present evidence, in this setting,
or provided a stated condition holds.
An \emph{enthymeme} leaves part of its reasoning unexpressed.

A \emph{critical question} probes a characteristic weakness of the selected
scheme. It may request support for a premise, test the relevance or strength
of an inference, suggest an exception, or introduce a competing argument.
An \emph{undercutting consideration} challenges the inferential connection
without necessarily supporting the contrary conclusion.
A \emph{rebutting consideration} supplies a reason for the contrary
conclusion. The term \emph{fallacious use} will identify a specified
misleading move, such as suppressing a relevant qualification or evading
a needed response; it will not label every instance of a defeasible scheme.
\end{hbdefinition}

\begin{hbmethod}{def:p18-scheme-reconstruction-method}{Reconstruct before assessing}
Use the text, its immediate context, the question being discussed, and
the available supporting records as inputs. Produce an explicit
reconstruction together with a qualified assessment.

\begin{enumerate}
\item Identify the intended conclusion and its role. Is the speaker
asserting a fact, proposing an explanation, recommending an action,
or assigning responsibility? Record the population, time, place,
conditions, and degree of confidence actually claimed.
\item Separate what the speaker states from what the context supplies.
Mark an added premise as an \emph{attributed implicit commitment}
only when textual or contextual evidence supports that attribution.
Otherwise call it a \emph{needed but unestablished premise}.
\item Treat scheme identification as an interpretive hypothesis.
Compare possible readings. A report of a perception may be witness
testimony; a conclusion drawn from that report may additionally require
classification, causal reasoning, or analogy.
\item Write the scheme's premises, warrant, and qualified conclusion for
this case. Preserve the distinction between a descriptive conclusion
and a recommendation. Do not make the reconstruction deductively valid
merely by adding the unsupported premise that its premises imply its
conclusion.
\item Ask the questions of the identified source version, one at a time.
For each question record its target, the available answer, any evidence
for that answer, and the remaining uncertainty. Distinguish asking
whether an exception exists from supplying grounds that it does.
\item Assess the argument at the claimed strength. Possible outcomes
include provisional support, a narrower supported conclusion, suspended
support pending an answer, a defeated inference, and an unsupported
claim. Explain which premise or inferential step determines the result.
\end{enumerate}

This is an interpretive and normative method. It is not a decision
procedure with a completeness theorem, and it supplies no empirical
accuracy rate. Its interpretation must remain answerable to the text;
its evaluation must remain answerable to the evidence.
\end{hbmethod}

Godden and Walton's reconstruction account treats schemes as models
against which an interpretation is checked, rather than as devices that
automatically supply whatever premises would rescue an argument
\cite[pp.~271--272]{GoddenWalton2007SchemesQuestions}.
If a proposed implicit premise conflicts with the speaker's stated
commitments, retain that conflict. It may indicate a different intended
scheme, an incomplete argument, or a misuse of the scheme.

Three standards must remain separate. Deductive validity excludes a case
in which the premises are true and the conclusion false; a sound
deductive argument also has true premises. Defeasible support can be
reasonable while leaving open exceptions and contrary evidence.
The adequacy of an answer in a dialogue additionally depends on what
was asked, what has been conceded, and which party must support the
relevant assertion. Godden and Walton assess cogency through premise
acceptability, relevance, and sufficiency; their scheme-specific questions
apply those standards to characteristic failure conditions
\cite[pp.~277--283]{GoddenWalton2007SchemesQuestions}.

\section{Argument from expert opinion}
\label{sec:p18-scheme-expert}

\begin{hbdefinition}{def:p18-expert-opinion-scheme}{Expert opinion: the selected presentation}
Let \(E\) be a source, \(S\) a subject domain, and \(A\) a proposition.
The selected scheme has the following content:
\begin{enumerate}
\item \(E\) has relevant expertise in \(S\), and \(A\) falls within \(S\).
\item \(E\) asserts \(A\), or asserts its negation.
\item Presumptively, the proposition asserted by \(E\) is true.
\end{enumerate}
The third item is the conclusion, interpreted defeasibly.
Walton's 2016 presentation compares a version with domain membership as
a separate premise, a version placing it in the expertise premise,
and an expanded conditional version. The expanded conditional is
defeasible; it does not declare experts infallible
\cite[Sections 4.4--4.6]{Walton2016ArgumentEvaluation}.
\end{hbdefinition}

The six basic questions are:

\begin{enumerate}
\item \emph{Expertise:} What supports the source's credibility as an expert?
\item \emph{Field:} Does the source's expertise cover the subject of this
particular proposition?
\item \emph{Opinion:} What exactly was asserted, and does its content
support the proposition attributed to the expert?
\item \emph{Trustworthiness:} Is the source personally reliable in the
present use of the testimony?
\item \emph{Consistency:} How does the proposition compare with the
judgments of other relevant experts?
\item \emph{Backup evidence:} Is the assertion evidence-based, and
what is that evidence?
\end{enumerate}

These are the six questions in the selected Godden--Walton and Walton
presentations, not an exhaustive account of all expert assessment
\cite[pp.~278--279]{GoddenWalton2007SchemesQuestions}
\cite[pp.~130--134]{Walton2016ArgumentEvaluation}.
Expertise and trustworthiness are distinct: a technically competent source
can be unreliable on an occasion, and an honest source can lack the
required expertise. Disagreement invites examination of methods and
evidence; counting expert names does not resolve it.

The earlier 1989 version reproduced by Walton denotes the subject domain
by \(D\) and has three separate premises: \(E\) is expert in \(D\);
\(E\) asserts that \(A\) is known
to be true; and \(A\) belongs to \(D\). Its conclusion is plausible
acceptance of \(A\). Its six questions concern, in order, fit to the
expert's competence, genuine expertise rather than celebrity, degree
of authority, disagreement by other experts, agreement with objective
evidence, and accurate interpretation of the pronouncement.
This is the older list, not another ordering of the later list.
The expanded conditional version adds the defeasible warrant from
relevant expertise and the assertion to the proposed acceptance of
\(A\); the qualification is part of how that warrant is used.

Walton also identifies three further questions under the expertise
heading: possible bias, honesty, and conscientiousness
\cite[pp.~129--132]{Walton2016ArgumentEvaluation}.
They show how one basic question can lead to more detailed examination,
and how question headings can overlap. They do not justify an endless
sequence of unsupported allegations.

\begin{hbexample}{def:p18-expert-opinion-worked}{An environmental-control recommendation}
Suppose an archive receives this statement:
\emph{Rao, an engineer who designs environmental controls for archives,
says the proposed controller kept Room 3 within the specified humidity
range during the six-week trial and is likely to continue doing so
provided the tested load and power conditions continue.
We should therefore treat that controller as a
promising option for Room 3.}

The supported factual conclusion is conditional and local: the controller
appears capable of maintaining the specified range under the trial's
conditions. The final recommendation also requires practical premises
about costs, alternatives, and the archive's needs.
The engineer's relevant competence and the actual assertion are stated.
That the report accurately represents the trial, that the instruments
were suitable, and that the proposed use sufficiently resembles the trial
are further commitments to investigate.

The expertise question calls for the engineer's work in this kind of
control system, not merely an impressive title. The field question
distinguishes engineering control of humidity from a claim that every
kind of object will thereby be preserved. The opinion question requires
the report and its qualification about loads and power.
If the trial was funded by the supplier, the trustworthiness question
asks about disclosure, reporting choices, and independent checks; funding
alone does not prove dishonesty. For consistency, an independent engineer's
objection about power interruptions should be examined on its merits.
For backup evidence, request the observations, instrument checks,
conditions, and basis for extrapolating beyond six weeks.

To finish the case, stipulate that the archive verifies Rao's work in
this field, checks the attribution against the signed report, and obtains
the trial observations and instrument records. The supplier's funding
is disclosed. A consulted independent engineer accepts the limited trial
finding but points out that no power-failure test was performed.
There is no evidence of fabricated or selectively withheld observations
in this stipulated record.
The appropriate outcome is provisional support
under the stated operating conditions, with the interruption question open.
It is neither proof of indefinite preservation nor a completed procurement
decision.

A misleading version says:
\emph{A leading expert guarantees that the whole collection is safe;
further questions only show distrust of expertise.}
This changes a bounded technical finding into an unrestricted guarantee,
crosses from control performance to preservation of every object, and
tries to prevent examination of the evidence. Correct the conclusion
and restore the omitted conditions. The correction does not establish
that the controller is ineffective; it establishes that this stronger
claim lacks the offered support.
\end{hbexample}

\subsection{The question's burden depends on the selected version}

Walton's 2016 Carneades representation has three ordinary premises:
the expert's domain, the proposition's membership in that domain,
and the expert's assertion. It treats evidential backing as an
assumption, and personal unreliability and conflicting expert opinion
as exceptions \cite[pp.~133--134]{Walton2016ArgumentEvaluation}.
On that particular assignment, a challenge to evidential backing puts
the initiative on the proponent to provide it. An allegation of
unreliability requires supporting evidence from the challenger before
it functions as the specified exception.

This representation is distinct from the preceding six-question display.
It supplies a rule for managing questions in that model; it is not a
universal rule that every accusation must be disproved by the expert,
or that every bare question already defeats an argument in every dialogue.
If a question introduces a new factual allegation, record that allegation
and its support instead of hiding it in interrogative wording.

\section{Argument from witness testimony}
\label{sec:p18-scheme-witness}

\begin{hbdefinition}{def:p18-witness-testimony-scheme}{Witness testimony: the three-premise version}
For a witness \(W\) and proposition \(A\), Walton's selected presentation
requires that \(W\) be in a position to know whether \(A\), that \(W\)
report truthfully according to what \(W\) knows, and that \(W\) state
\(A\) or its negation. The stated proposition may then be accepted
plausibly, subject to further assessment
\cite[pp.~271--273]{Walton2016ArgumentEvaluation}.

Being in a position to know includes an appropriate opportunity to
observe or otherwise acquire the information. Truthful reporting is
not the same as an accurate memory or an accurate interpretation of
what was observed.
\end{hbdefinition}

The five questions in this presentation are:

\begin{enumerate}
\item Is the witness's account consistent within itself?
\item Does it agree with established case information obtained
independently of this testimony?
\item Does it agree with other witnesses' independent testimony?
\item What grounds, if any, support attributing bias to this account?
\item How plausible is the particular proposition being asserted?
\end{enumerate}

The position-to-know and truth-telling premises remain part of the
scheme even though the five-question list does not repeat them word
for word. Questions about visibility, memory, and reporting conditions
may examine those premises. Do not treat several retellings of one
person's story as independent corroboration.

\begin{hbexample}{def:p18-witness-testimony-worked}{What the gate witness did and did not see}
Suppose Sam states:
\emph{At about 6:10 I was closing the east service gate. I saw a yellow
van parked across the road. I did not see who was inside.}
The context supplies a signed report made that evening.
A defensible reconstruction concludes that there is testimonial support
for a yellow van being at that location around that time.
It also makes visible the assumptions that Sam could see the van,
remembered the occasion, and reported the experience sincerely.

For internal consistency, compare the time and location throughout
Sam's account. A discrepancy between 6:10 and 6:30 needs an explanation,
but could concern a clock or an estimate rather than fabrication.
For consistency with other evidence, compare the gate log and any
independent recording, including their own time and location uncertainties.
For independent testimony, distinguish a second observer from a colleague
who merely repeats Sam's report. For bias, investigate a concrete motive
or influence where there is evidence for it; an unsupported insinuation
is not a finding. For plausibility, ask whether the stated observation
fits the viewing conditions and ordinary circumstances of the case.

To finish the case, stipulate that the evening report is internally
consistent, the gate record fixes Sam's approximate location and time,
and a second observer independently reports the same van.
The viewing conditions permitted that observation. The supplied record
contains no evidence of a motive to misdescribe the van.
Neither observer identifies its occupant. The conclusion
remains about the van. Further inferences about ownership, occupancy,
and a connection with missing property need their own premises.

The misleading reconstruction is:
\emph{Sam saw Alex steal the missing item.}
It substitutes an inferred identity and action for an expressly limited
observation. Even if Alex owns a yellow van, ownership would not by itself
identify the driver on this occasion. Repair the report to preserve what
was actually witnessed, and separately assess each proposed connecting
inference. Failure of the theft inference does not show that Sam's
description of the van was false.
\end{hbexample}

This entry assesses evidential reasoning, not the current admissibility
rules of any court. A formal legal burden or exclusion rule would require
an identified jurisdiction and its own source. The witness questions
also do not supply a general psychological accuracy rate.

\section{Argument from analogy}
\label{sec:p18-scheme-analogy}

\begin{hbdefinition}{def:p18-analogy-basic-scheme}{Source and target cases}
Walton's 2014 basic presentation distinguishes a source case \(C_1\)
from a target case \(C_2\). A situation is described in \(C_1\);
a conclusion \(A\) is plausibly drawn from that case; and \(C_1,C_2\)
are claimed to be similar. The argument transfers a comparable,
plausibly acceptable conclusion to \(C_2\)
\cite[author version, pp.~1--3]{Walton2014AnalogySchemes}.

The earlier compact form cited there says that \(A\) holds in \(C_1\)
and transfers it to a similar \(C_2\). The revised form makes the
source-case inference itself explicit. Consequently the source outcome
is assessable; it is not automatically a sound premise just because
the story is familiar or rhetorically effective.
\end{hbdefinition}

The three matching questions are:

\begin{enumerate}
\item Are there differences between \(C_1\) and \(C_2\) that weaken
the particular similarity being relied upon?
\item Is \(A\) an appropriate conclusion in the source case \(C_1\)?
\item Is there another case \(C_3\), similar to \(C_1\), in which
a different conclusion should be reached?
\end{enumerate}

The third is the source presentation's counteranalogy question. A critic
must explain the relevant comparison and the different warranted outcome;
merely telling a different story is not yet a counterargument.
Likewise, a difference matters only through its bearing on the issue.
Two cases can share many superficial features while differing on the
feature that explains the proposed treatment.

\subsection{The distinct factor-based presentation}

Walton also presents a second approach that specifies the features
shared by the objects or cases, rather than leaving similarity
unanalyzed. In the core form discussed through Guarini, both cases
have the listed features and are proposed for the same treatment or
classification with respect to those features.
The assessment then asks which similarities support the treatment
and which disanalogies oppose it
\cite[author version, pp.~5--8]{Walton2014AnalogySchemes}.

The textbook forms compared in that discussion should also be kept
distinct. In the Copi--Cohen form, four entities \(a,b,c,d\) share
attributes \(P,Q\), the first three also have \(R\), and \(d\)
is inferred probably to have \(R\). In the Hurley form, one entity
has the listed shared attributes and an additional attribute, and
that additional attribute is tentatively transferred to a second entity.
These are attribute-transfer patterns, whereas the core feature-based
scheme proposes common treatment or classification.
Walton reports Hurley's six appraisal criteria: relevance and number
of similarities; the nature and degree of differences; the number
and diversity of the source analogues; and the specificity of the
conclusion. The last criterion makes the conclusion's strength an
explicit part of the assessment.

This is not a license to transfer every attribute of one object to
another. Neither is it a rule that the larger count of favorable
features wins. Relevance to the issue and the significance of each
feature need reasons. Walton's three-move comparison starts with
a supporting case, permits a distinguishing factor or counteranalogy,
and then permits a response addressing that challenge.
The source's later script approach additionally compares organized
sequences of actions and circumstances, rather than isolated matching
attributes \cite[author version, pp.~11--15]{Walton2014AnalogySchemes}.

\begin{hbexample}{def:p18-analogy-worked}{A late-return waiver}
Suppose a department waived a laptop late fee when a university server
outage prevented a borrower from renewing an otherwise renewable loan.
A second borrower asks for a camera-loan fee to be waived:
\emph{The same outage stopped my renewal, so the department should
waive this fee too.}

The source case is the laptop waiver; the target is the camera loan.
The intended conclusion is equal treatment for the delay attributable
to the outage, not abolition of all late fees.
The outage and source waiver are stated. Potential implicit commitments
include that renewal would otherwise have been permitted, that the
borrower attempted it in time, and that the reason for the source waiver
applies to camera loans. These commitments must be checked rather
than supplied merely to make the analogy persuasive.

For the first question, a relevant difference would be a reservation
that made this camera nonrenewable even without the outage.
The different price or color of the devices matters only if connected
to the waiver policy's purpose.
For the second question, examine why the laptop waiver was justified:
perhaps the department accepted responsibility for its inaccessible
renewal service; perhaps instead it was an unrelated discretionary favor.
The latter would weaken the proposed generalization.
For the third question, consider a borrower whose loan could not be
renewed because another user had reserved it, although the same outage
occurred. If charging that borrower was appropriate, the source and
counterexample suggest that technical inability to renew is relevant
only where renewal would otherwise have been available.

The factor-based assessment therefore distinguishes the common outage,
the attempted renewal, renewability, and the reason for the waiver.
In the completed example, the renewal record shows a timely attempt,
the camera had no reservation restriction, and the laptop decision
expressly relied on the failure of the renewal service.
For this assessment, both parties accept the department's stated
principle that a borrower should not pay for a delay caused by that
failure when renewal was otherwise available.
The analogy therefore gives a defeasible reason for a comparable waiver.
If the reservation restriction applies, this renewal-based analogy
does not justify the waiver; a different argument would be needed.

A misleading appeal says:
\emph{Both are university devices, so any different decision is
inconsistent and unfair.}
It suppresses the possible difference that explains the treatment.
The correction is to identify the operative reason for the source
decision and compare the target on that reason. Merely alleging
inconsistency does not establish that the cases deserve identical results.
\end{hbexample}

\section{Causal inference, appearances, and signs}
\label{sec:p18-scheme-causal}

\begin{hbdefinition}{def:p18-correlation-cause-scheme}{A correlation as a defeasible reason for a causal hypothesis}
The correlation-to-cause scheme takes a positive association between
events or kinds of events \(E_1,E_2\) as a defeasible reason for
the hypothesis that \(E_1\) causes \(E_2\).
An association requires an appropriate comparison of occurrences or
measurements; the mere fact that two events occurred is not a
statistical finding of association.

Walton's selected treatment distinguishes this provisional inference
from a causal conclusion adequately supported for a later stage of
an inquiry. A reported association is not itself an experimental
demonstration of an intervention's effect
\cite[Sections 6.2 and 6.9--6.10]{Walton2016ArgumentEvaluation}.
\end{hbdefinition}

The 2016 revised list has ten questions. Its first seven preserve
the earlier list discussed in Section 6.2; the last three add a
theoretical connection and the distinction between a surrogate
measurement and the outcome actually at issue:

\begin{enumerate}
\item Is the asserted positive association present in the observations?
\item Are there enough instances to support the claimed association?
\item What evidence favors the direction from \(E_1\) to \(E_2\)
over the reverse direction?
\item Could a third factor account for both \(E_1\) and \(E_2\)?
\item If intermediate variables occur, is the proposed causal connection
indirect, through those variables?
\item What range of cases supports the association, and where does
it cease to hold?
\item Could the apparent change in \(E_2\) be produced by a changed
definition, classification, or reporting standard?
\item Is there a scientific theoretical basis for the proposed
connection between \(E_1\) and \(E_2\)?
\item Is the argument about the final outcome of interest or about
a surrogate used in its place?
\item If a surrogate is used, what justifies the step from that surrogate
to the final outcome?
\end{enumerate}

The list is not a scoring rule. Ten favorable answers do not yield
a numerical probability or eliminate every possible alternative.
The mediation question also differs from the common-cause question:
an intermediate stage in the proposed mechanism is not automatically
an alternative explanation that removes that mechanism.

\begin{hbexample}{def:p18-causal-worked}{A rush-processing procedure and damaged shipments}
Suppose a manager reports:
\emph{There were 12 damage reports per 100 shipments using the new
rush procedure and 4 per 100 shipments using the ordinary procedure.
The new procedure therefore causes three times as much damage and
must be banned at every depot.}
Treat these numbers as stipulated case data, not research findings.

The records describe reports, whereas the factual conclusion concerns
damage to the shipped contents. The proposed ban is a further practical
conclusion. Reconstruct these as separate steps. The reporting association
is stated. The claim that shipments were comparable, that reporting was
consistent, that the direction is causal, and that the measured report
category represents actual damage are needed additional premises.

The first question checks the recorded numerators, denominators, and
comparison groups. The second asks whether the observations cover enough
independent circumstances, or mostly one unusually difficult delivery run.
The third asks whether the rush procedure was introduced because certain
routes already had damage problems. The fourth asks whether fragile
contents, a particular carrier, or workload differences explain both
procedure use and the reports.
The fifth examines an actual proposed sequence, such as reduced packing
time leading to omitted padding and then breakage.
The sixth restricts any conclusion to the types of goods, workers,
routes, and conditions that the evidence can support.

For the seventh question, suppose the rush team began recording cosmetic
carton dents that the ordinary team had not recorded. That changes the
measurement. For the eighth, a specific account of how the changed
handling could produce damage needs investigation.
The ninth asks whether a damaged carton is being used as a surrogate
for broken contents. The tenth asks whether inspection of the contents
supports that inference. A dent can be real even when the contents
are intact; the two propositions must not be substituted for one another.

To finish the case, stipulate that rush shipments disproportionately
contained fragile goods, that the new reporting category was used only
by the rush team, and that no independent inspection of contents was
performed. These give an actual comparability problem, an actual change
of classification, and an unverified step from carton reports to contents
damage. The timing of procedure adoption and the proposed mechanism
remain uninvestigated in the supplied record.
The comparison therefore does not establish the manager's causal conclusion.
The data can still motivate a focused investigation.
Comparable measurements, appropriate comparisons, and evidence bearing
on the mechanism may strengthen or weaken the hypothesis.
The argument for a company-wide ban additionally requires feasible
alternatives, costs, competing goals, and the expected consequences
of that action.

The misleading move is not the act of noticing a correlation.
It is the jump from an inadequately examined comparison to a quantified
causal conclusion and an unrestricted policy, while suppressing the
questions that bear on that jump. The repair is a properly qualified
claim and an explicit investigation of the missing connections.
It does not establish that the procedure has no harmful effect.
\end{hbexample}

\subsection{The source's distinct burden representation}

The Carneades representation discussed by Walton contains an ordinary
premise that the events are correlated, an assumption that a theory
explains the proposed connection, and an exception specifying a common
cause. On that representation, asking for the explanatory basis requires
the proponent to supply it; a common-cause exception requires the
challenger to identify and support the proposed factor
\cite[pp.~185--186, 203--204]{Walton2016ArgumentEvaluation}.
This smaller representation does not replace the ten-question inquiry
list. Nor does it make the absence of a presently identified common
cause a deductive proof of causation.

\subsection{Forward causal reasoning and backward signs}

Macagno, Walton, and Reed distinguish the causal relation from the
direction of reasoning that uses it. For event descriptions \(A,B\),
their cause-to-effect pattern combines a defeasible generalization
that \(A\) normally produces \(B\)
with information that \(A\) occurs, and tentatively concludes \(B\).
Reasoning backward from an observed \(B\) to a possible \(A\) is a
different inferential step
\cite[pp.~2495--2496]{MacagnoWaltonReed2017Schemes}.

For example, a wet patch can be treated as a sign of a leak only through
some connection between the observation and that proposed explanation.
It could also be explained by cleaning, a spill, or a mistaken
classification of a reflection as water. Observing the patch does
not by itself establish a unique cause.

Walton separately gives an appearance-based scheme: an image makes
an object appear to fall within a category, so the object is tentatively
classified that way. His two questions ask whether that appearance
could mislead and whether another classification has better support
\cite[pp.~27--31]{Walton2016ArgumentEvaluation}.
Those are questions for that appearance presentation; they are not
declared a complete question set for every argument from sign.
First assess whether the surface is actually wet, then assess what
would explain its being wet. An accurate observation can coexist
with an unsuccessful explanation.

\section{Practical reasoning from goals to action}
\label{sec:p18-scheme-practical}

\begin{hbdefinition}{def:p18-practical-reasoning-scheme}{The Godden--Walton necessary-condition presentation}
Let \(A\) denote a desired state and let \(B_0,\ldots,B_n\) be a finite
nonempty list of considered actions, with \(n\geq0\).
The selected action has index \(i\in\{0,\ldots,n\}\).
The selected presentation has five premise roles:
\begin{enumerate}
\item The agent aims to bring about \(A\).
\item On the available information, achieving \(A\) requires taking
at least one action from the considered alternatives.
\item The agent selects \(B_i\) as an acceptable, or the most acceptable,
necessary option.
\item No known unchangeable circumstance prevents the agent from
carrying out \(B_i\).
\item Achieving \(A\) is more acceptable to the agent than refraining
from \(B_i\), with its consequences taken into account.
\end{enumerate}
The conclusion is that the agent is required to carry out \(B_i\),
understood as a defeasible practical requirement relative to the goal,
the alternatives, and those assessments
\cite[pp.~285--287]{GoddenWalton2007SchemesQuestions}.
This is neither an unconditional moral obligation nor a deductive
consequence of the bare fact that an action could help achieve a goal.
\end{hbdefinition}

The five matching questions are:

\begin{enumerate}
\item Are there other available means of attaining \(A\)?
\item Is the proposed action acceptable, or preferable among the
relevant options?
\item Can this agent actually perform it?
\item What adverse side effects, if any, should enter the assessment?
\item What other goals might conflict with pursuing \(A\) in this way?
\end{enumerate}

An option that is sufficient for a goal is not thereby necessary,
and a necessary condition need not suffice by itself.
Selecting an option does not make the selection well supported.
The goal, comparison of alternatives, feasibility, and consequences
must all remain open to assessment.

\begin{hbexample}{def:p18-practical-reasoning-worked}{Catalogue access during a server replacement}
Suppose a library plans a week of server replacement and a manager says:
\emph{Readers need catalogue searches during that week. We should
provide a temporary read-only copy.}
The stated goal is continued search access during a defined interval.
The proposed action is a temporary copy, not permanent replacement
of every library system.

The reconstruction needs a premise about the considered options:
a read-only copy, staff searches using an existing printed index,
and access through a partner service, for example.
It needs reasons that some available measure is required during the
outage, that the selected copy is an acceptable choice, that the
library can prepare and run it, and that its benefits outweigh its
disadvantages relative to the other goals.
None of those claims follows from the first sentence alone.

The alternatives question asks about existing services and different
maintenance scheduling. The selection question compares the useful
coverage, effort, and cost of the realistic alternatives.
The possibility question checks the available export, equipment,
permissions, and staff time. The side-effects question examines stale
records, confusing availability information, and additional work.
The conflicting-goals question considers commitments such as accurate
information and protecting access-restricted records.
These case-specific investigations instantiate the five questions;
they do not add five assumed facts.

In the completed case, the printed index is three months out of date,
and the partner service does not cover the archive's finding aids.
A current export of the public catalogue is available and authorized;
it excludes access-restricted records. An existing machine and an
identified half-day of staff time suffice for the specified setup.
The copy will clearly state that live availability is unavailable.
On these stipulated facts the copy better meets the temporary search
goal than the two considered alternatives, without taking on the
obligations of a permanent replacement.
That supplies a defeasible, bounded reason to proceed.
If those facts are absent, the conclusion remains a proposal awaiting
support rather than an established requirement.

A misleading version says:
\emph{Catalogue access is essential, so we must immediately purchase
the proposed permanent platform.}
It substitutes a different action, omits alternatives, and treats one
possible means as necessary. The repair is to restore the temporary
goal, compare available means, and justify the selected action.
Acceptance of the goal does not compel acceptance of that purchase.
\end{hbexample}

\subsection{Probing an option and asserting an objection}

Godden and Walton distinguish a question used to probe the option search
from one used to assert that the search missed a better option.
The first invites the proponent to explain what was considered.
The second incurs a burden to support the alleged alternative
\cite[pp.~283--287]{GoddenWalton2007SchemesQuestions}.

In the library case, \emph{What alternatives did you consider?} can
receive a reasoned account of the options already investigated.
\emph{Why ignore the partner service we can already use?} additionally
asserts that this service is available and relevant. Its proponent
must support those claims; the manager must then address their bearing
on the choice.
Repeatedly raising an unspecified possibility is not the same as
establishing an overlooked alternative. Conversely, a merely perfunctory
answer does not settle whether a concrete, evidenced objection has
been answered. The source's burden discussion is a proposed account
of these dialogical roles, not a universal rule for every institutional
decision procedure.

\section{From a question list to an assessment}

The appropriate final record names the conclusion assessed, the supporting
scheme, the evidence for each needed premise, the critical questions
actually raised, the answers, and the resulting qualification.
It distinguishes a supported objection from a possible line of inquiry.
It also records what would change the assessment.

For the cases above, the distinction is consequential. A conditional
expert judgment does not authorize a universal guarantee. A witness's
observation does not automatically identify an actor or establish an act.
An analogy requires a defensible source conclusion and relevant similarity.
A correlation does not by itself establish a causal effect or an
appropriate intervention. A goal does not determine its preferred means.
In each case the defect can be located and the conclusion corrected
without asserting that every argument of that type is fallacious.

Godden and Walton's account also explains why asking every listed
question is not a mechanical completion criterion. Their questions
address characteristic ways of failing acceptability, relevance, or
sufficiency; additional circumstances can matter, and a later discovery
can change the assessment. Whether a question shifts a burden depends
on its role and on the dialogue's commitments
\cite[pp.~280--287]{GoddenWalton2007SchemesQuestions}.

The structured framework in
Chapter~\ref{sec:p18-structured-argumentation} can represent an explicitly
formulated premise challenge, undercutter, or rebuttal.
An interrogative sentence alone is not an argument tree with premises
and a conclusion. To represent a critical response formally, specify
what it asserts, what supports that assertion, and which target it
challenges. The resulting formal model supplements the reconstruction;
it does not supply an omitted observation, establish an expert's
competence, or decide the relevance of an analogy by itself.

\begin{proofboundary}{def:p18-schemes-next-boundary}{Further source presentations}
These entries cover the identified expert, witness, analogy, causal,
appearance, and practical-reasoning presentations. The complete 2017
classification, its historical comparisons, and worked diagnosis of
connected schemes are supplied in
\cref{sec:p18-scheme-classification-diagnosis}. The complete historical
scheme catalogues, other testimony and practical-reasoning variants,
the original historical classification works, and the separate argument-from-sign
catalogue entry remain distinct source obligations.
Broader treatments of informal fallacies, rhetoric, relevance, dialogue
rules, and institutional burdens also require their own presentations.
No completeness theorem, general empirical effectiveness claim, or
universal burden protocol is asserted for this method.
\end{proofboundary}
